\section{A smaller arithmetic certificate base: \texorpdfstring{$F=t^t$}{the diagonal power}}\label{sec:selfpower}

This section belongs to the direct proof-calculus construction and uses a
different $F$ from the compact mask. The factorial-free predecessor uses
$F=[2(t+1)]^t$ to separate proof-label multiplicities and make its binomial
CRT modulus cover every row. The diagonal power $F=t^t$ suffices for both
jobs. Its exact power relation is compiled into polynomial Pell residuals;
it is not left in the final Diophantine equation.

\begin{lemma}[Label balance at the diagonal power]\label{lem:labels}
Let $t\ge1$, $F=t^t$, and $W_j=1+jF$ for $0\le j\le t$.
Suppose $v_i,p_i,q_i\in\{0,\ldots,t\}$ for $0\le i<t$, and
\begin{equation}\label{eq:labels}
 \prod_{i<t}W_{v_i}=(1+F)\prod_{i<t}W_{p_i}W_{q_i}.
\end{equation}
At every positive label $j$, its conclusion multiplicity minus its
premise multiplicity is one at $j=1$ and zero otherwise.
\end{lemma}
\begin{proof}
Discard the label-zero factors, which are one. If the two factor counts
differ, their smaller count $r$ is at most $t$, because the left side
has at most $t$ nonunit factors. The side with more factors is at least
$(F+1)^{r+1}$, and the other side is at most
$(1+tF)^r\le t^r(F+1)^r$. But $F+1=t^t+1>t^r$, a contradiction.
Hence their counts agree, with common value $r\le t$.

Replace $F$ by an indeterminate $z$. Each side is a product of $r$
linear factors $1+jz$, where $1\le j\le t$. Its coefficient of $z^k$,
for $0\le k\le r$, is positive and at most
\[
 \binom rk t^k\le t^{r-k}t^k=t^r\le F.
\]
For the inequality, inject the $(r-k)$-element subsets of an $r$-element
set into ordered $(r-k)$-tuples from a $t$-element set. Thus every
coefficient of the difference polynomial has absolute value strictly
less than $F$. If the difference were nonzero, divide its evaluation
at $F$ by the power corresponding to its lowest nonzero coefficient
and reduce modulo $F$. That coefficient would be a nonzero multiple
of $F$ of absolute value less than $F$. This is impossible.
The products are therefore equal in $\Z[z]$, and the distinct roots
$-1/j$ give the claimed multiplicities. When $F=1$, the coefficient
difference bound itself already implies that the difference is zero.
\end{proof}

The strict bound is on the \emph{difference} coefficients. Bounding
each entire product by its coefficient sum misses this improvement
and retains the unnecessarily larger base $(t+1)^t$.

\begin{lemma}[Single-conclusion flow]
Take a finite multiset of valid Horn rules with one conclusion and at most
two premises; pad unused premises with freely derivable truth, label zero.
If their multiplicities have the balance in Lemma~\ref{lem:labels},
then label one has a finite derivation.
\end{lemma}
\begin{proof}
Form the closure of truth under these finitely many rule instances.
If label one were outside, every rule whose conclusion is outside would
have at least one premise outside. The sum of conclusions minus premises
over outside labels would be at most zero. The prescribed balance
instead gives one, a contradiction.
\end{proof}

This inherited flow argument explains why an unordered certificate with
cycles does not create a false proof. It needs no ordering of the rows.

\subsection{The binomial compiler still covers every row}

Let $\Lambda_t=\lcm(1,\ldots,t)$ and $\mathcal R=F^2$. We need only
\emph{numerical domination}, not $\Lambda_t\mid F$:
\[
 \Lambda_t\le t!\le t^t=F,\qquad \mathcal R\ge\Lambda_t^2.
\]
Indeed $\Lambda_t$ need not divide $t^t$: three divides $\Lambda_{16}$
but not $16^{16}$. Avoiding this divisibility requirement was the
purpose of the inherited factorial-free compiler.

For clarity, its two prime-power facts are recalled here. Put
$\mathcal B=\binom{\mathcal R+t}{t}$. Every full prime power
$p^e\Vert\mathcal B$ divides some $\mathcal R+j$, $1\le j\le t$:
Legendre's valuation formula consists of zero-one summands and its
sum is at most $\max_jv_p(\mathcal R+j)$. Conversely, each row $j$
has a full prime power of $\mathcal B$ exceeding $t$. Since
$\mathcal R+j>\Lambda_t^2$, some $v=v_p(\mathcal R+j)$ exceeds
$2\lfloor\log_p t\rfloor$. For this prime, the valuation in
$\mathcal B$ is exactly
\[
 v-v_p\!\left(t\binom{t-1}{j-1}\right).
\]
The subtracted valuation is at most $\lfloor\log_p t\rfloor$, because
$t\binom{t-1}{j-1}\mid\Lambda_t$. This classical denominator bound
also follows from Farhi's binomial lcm identity~\cite{farhi}.
The resulting full prime power therefore exceeds $t$ and identifies
the row uniquely. The first fact covers \emph{every prime} in the
converse CRT construction, while the second covers \emph{every row}
in soundness. Both remain valid with $F=t^t$.

\subsection{Parity and cofinal completeness}

The retained horizon is $t=16+h_0^2$. For completeness choose the
cofinal subfamily with even $h_0$, so $t$ is even. Then
\[
 v_2(F)=t\,v_2(t)\ge t.
\]
Thus $F$ and the beta base $y=F(2u+1)^2$ are even.
Let $O$ be the odd part of $t!$. Choose $2u+1$ to be a sufficiently
large positive odd multiple of $O$. Since $v_2(F)\ge t>v_2(t!)$,
this makes $y$ an even multiple of $t!$, as required by the inherited
pairwise-coprime beta-modulus construction. This choice supplies the
missing odd prime factors; it does not incorrectly assume $t!\mid t^t$.
The beta moduli $1+(j+1)y$ are all odd, so their CRT progression
admits arbitrarily large even representatives for the beta code.
The binomial modulus $\mathcal B$ is odd: since
$v_2(F^2)>\lfloor\log_2t\rfloor$,
$v_2(F^2+j)=v_2(j)$ for $1\le j\le t$, and the numerator and $t!$
have equal two-adic valuations. All inherited CRT parity choices
remain available.

Soundness imposes no new evenness condition on an arbitrary satisfying
horizon. The independent audit examines each use of the old base
$[2(t+1)]^t$, rather than replacing it in the formula and assuming
all estimates survive. The needed size bound follows already from
$F=t^t\ge16^{16}=2^{64}$ at $t\ge16$, together with
$F\ge\Lambda_t$. The source edit also changes the actual power-call
base and the shared frame metadata, which are checked separately.

\section{Eliminating common quadratic directions}\label{sec:radical}

The second direct improvement is algebraic and reusable. An ordinary
monic substitution can raise degree: replacing a variable by a quadratic
expression where it appears quadratically elsewhere may create quartic
residuals and an eighth-degree sum of squares. Instead, we find directions
along which \emph{every} residual's quadratic part is constant.

Write
\[
 R_j(\boldsymbol x)=c_j+\ell_j(\boldsymbol x)+Q_j(\boldsymbol x),
\]
where $\ell_j$ and $Q_j$ are homogeneous of degrees one and two.
Let $H_j$ be the integer Hessian matrix of $Q_j$; a term $a x_r^2$
contributes $2a$ in its diagonal position.

\begin{definition}
The common quadratic radical is
\[
 \mathcal K=\bigcap_j\ker_{\Q}H_j\subseteq\Q^n.
\]
\end{definition}

\begin{theorem}[Integral radical projection]\label{thm:radical}
Suppose $\boldsymbol d\in\Z^n\cap\mathcal K$ has $d_v=1$ at a
coordinate $v$, and $\ell_r(\boldsymbol d)=\varepsilon\in\{1,-1\}$
for a residual $r$. The common integer zero set of the system is then
in explicit polynomial bijection with the common integer zero set
of $m-1$ quadratic residuals in $n-1$ coordinates.
\end{theorem}
\begin{proof}
Make the integral coordinate change
\[
 x_v=y_v,\qquad x_w=y_w+d_wy_v\quad(w\ne v).
\]
Its inverse is $y_v=x_v$, $y_w=x_w-d_wx_v$, so it is unimodular.
Since $H_j\boldsymbol d=0$, the quadratic part of every transformed
residual is independent of $y_v$. Each transformed residual is
\[
 S_j(\widehat{\boldsymbol y})+\ell_j(\boldsymbol d)y_v,
\]
where the hat omits coordinate $v$ and $S_j$ is quadratic.
The selected row has the unique integral solution
$y_v=-\varepsilon S_r$. Substituting in all other rows gives
\[
 S_j-\varepsilon\ell_j(\boldsymbol d)S_r,
\]
still quadratic. Delete the selected row and coordinate.
Restriction maps an old zero to a new zero. Restoring $y_v$ by
the displayed expression and undoing the shear gives its unique
inverse. All maps have integer coefficients.
\end{proof}

The sum of squares still has degree at most four. Its polynomial need not equal the
old polynomial; the zero sets correspond. This distinction matters
when comparing collected term counts and coefficient heights.

\subsection{A small example}

For
\[
 R_1=(x+y)^2+z-1,\qquad R_2=(x+y)^2+x+2y,
\]
the direction $(1,-1,0)$ belongs to both Hessian kernels. Substituting
$x=u$, $y=v-u$ gives $R_1=v^2+z-1$ and $R_2=v^2+2v-u$.
The second row restores $u=v^2+2v$, leaving only $v^2+z-1=0$.
Both the shear and the restoration are integral. The large certificates
use exactly this mechanism.

\subsection{The six actual projections}

At every stage, the supplied certificate records the ordered variable
names, the integer radical direction, the unit coordinate, the selected
row, the shears, the quadratic restoration, and the retained-coordinate
map. Each branch uses six such projections. Each step deletes exactly
its selected coordinate and row; no further unused variable is discarded
at this stage.

The independent checker verifies $H_j\boldsymbol d=0$ for every row,
the unit slope, both coordinate identities, the restored pivot equation,
and every substituted residual as an exact polynomial identity.
It expands and collects the final sum of squares. Earlier preprocessing
may remove unused coordinates; that separate existence-of-zeros step
is not described as a bijection of all raw assignments.

The search computes rational kernels and clears denominators, but its
output certificate is integral. No search heuristic is a mathematical
premise. Exact modular rank certificates show that the final stacked
Hessian matrix has full column rank modulo $101$ in all three branches.
The corresponding integer minor is nonzero, so the common radical is
zero over $\Q$. There are no further directions of this particular type
at these charts. This is not a global minimality statement.

\section{From proof calculus to the fixed integer quartics}

The direct branch uses the numerical proof calculus $C_{222}$. The
retained source compiles its finite numerical Horn rule families, an
unordered flow certificate, bounded range assertions, binomial CRT
row selection, and a synchronized Pell power representation. The
result is a system of quadratic residuals with fixed integer coefficients.
It has no Turing machine coefficient tuple and does not invoke Jones's
universal polynomial.

The $235$- and $236$-variable branches begin with the factorial-free
compiler~\cite{factorialfree}, change the power request to $t^t$ by
Section~\ref{sec:selfpower}, and apply the six radical projections.
Their support tradeoff reflects an earlier dense affine elimination
used only in the $235$ branch. The $234$-variable branch applies the
six projections to the shared-factorial-error quartic~\cite{sharederror};
it does not claim to remove that branch's factorial identity.

The relevant chain, formalizable internally in ZFC, is
\[
 \begin{gathered}
  \text{an integer zero of the final quartic}\\
  \Longleftrightarrow\text{a valid compiled arithmetic certificate}\\
  \Longleftrightarrow\text{a finite }C_{222}\text{ refutation}\\
  \Longleftrightarrow\neg\Con(\mathrm{ZF})
  \Longleftrightarrow\neg\Con(\mathrm{ZFC}).
 \end{gathered}
\]
The last step is syntactic equiconsistency, using the retained
constructible-model calibration. It does not assume a transitive
model or an externally standard model. The archived calibration
review identifies the source interfaces on both sides.

An external proof-assistant theorem about standard numerical objects
and an encoded object-language ZFC derivation are different artifacts.
This release supplies ordinary arguments formalizable in ZFC and
retains the inherited formal-interface sources. It does not claim
that every stage has been newly elaborated as a Lean, Rocq, or
object-ZFC proof term. A finite Python replay authenticates coefficients
and transformations; it does not replace unbounded Pell, CRT, and
syntactic calibration arguments.
